% Discussion --- plan \S4.15 (Phase 5).
% Sources: new_notes3 Part G \S32 (assay predictions), Part J \S39 (open
% tasks); new_notes Part II \S18 (confrontation with single-cell data),
% Part III (open tasks).

\section{Discussion}
\label{sec:discussion}

The chapter set out to complete the single-cell theory of the
birth--death--catastrophe process and to scale it up into a population
model forced by the intracellular dynamics. Both are done: the single-cell
description is now complete in distribution --- geometric states at every
time, a geometric quasi-stationary limit, and a burst-size law identical to
it --- and the population model exists, verified, and carries exact
consequences for fitting, for thresholds, and for the evolutionary question
of bursting versus budding. What remains is to say what the results imply
for measurement, how they connect to the chapters that follow, and what is
still open.

\subsection{Single-cell assay predictions for lytic systems}
\label{sec:discussion-assay}

For a genuinely lytic system --- \yp{}, \ft{}, lytic phage --- a
limiting-dilution or Hataye-type assay (one infected cell per well, counts
followed in time) measures the single-cell process directly, with no
population model interposed. The theory makes five predictions, all of them
parameter-free given $(\beta,\mu,\delta)$:
\begin{enumerate}
\item a fraction $b$ of wells show no release ever --- the mass of internal
extinction;
\item the cumulative release in a well is a step function: zero, then a
single terminal jump;
\item the jump times follow the density $\delta J(t)/(1-b)$;
\item the jump sizes follow the geometric law with ratio $1/a$ --- the
quasi-stationary distribution, by Corollary~\ref{cor:burst-qsd};
\item the between-well variance of total release is $\mathrm{Var}(W_\infty)$
of \eqref{eq:EW2}, a quantity no deterministic constant-rate model can
produce at all.
\end{enumerate}
The honest caveat of \S\ref{sec:hiv} applies: these predictions are too
strong for HIV as literally stated --- intermittent budding is not a single
terminal burst --- and they should be used for lytic hosts, or as the
bursting endpoint of the spectrum of \S\ref{sec:spectrum}.

\subsection{Consequences for fitting practice}

Three consequences for the way within-host models are fitted to data.
First, $p$ is not a constant of the pathogen: the effective release rate
$p_{\rm eff}(r)$ falls monotonically with the epidemic growth rate, from
$V_\infty/\mean{T_{\rm prod}}$ in the stationary phase towards $\delta$ in
the acute phase, and fits to acute and set-point data should therefore
yield systematically different estimates, in a direction and by an amount
the theory specifies (\S\ref{sec:peff-section}). Second, the eclipse-like
delay usually inserted into within-host models by hand emerges from the
intracellular dynamics with no extra parameter; its duration is the mean
productive lifetime's shadow, $\beta^{-1}\log\langle X\rangle_{\rm QS}$.
Third, the identifiability negative result of
\S\ref{sec:identifiability} says that no fit of $(p,d_{\Icell})$ --- from
any data, at any growth rate --- can recover the three intracellular rates;
a third observable is required, and the theory names the three cheapest
ones (burst-size dispersion, the no-burst fraction $b$, the shape of the
release-time density). An estimate of $p$ alone tells one almost nothing
about the cell; an estimate of $p$ plus any one of the three tells one
something testable.

\subsection{Connections forward}

The two-type process of Chapter~5 supplies the most immediate continuation.
In the GATE regime $\delta_1=0<\delta_2$ --- no burst before conversion ---
the no-catastrophe probability has $S'(0)=0$: a zero initial hazard, then
rising. That is precisely the eclipse-phase signature which within-host
models insert by hand, emerging mechanistically, with the eclipse duration
set by the conversion rate $\nu$. The renewal construction of this chapter
is ready to receive the two-type kernels; what blocks the way is the
two-type calculation itself (below). The budding--bursting comparison of
\S\ref{sec:flooding} finds its evolutionary home in Chapter~6, where the
spectrum of \S\ref{sec:spectrum} --- partial release, reset catastrophe,
self-excitation --- provides the intermediate phenotypes a selection
argument needs; and the nonlinear mechanisms Chapter~7 discusses in the
Komarova vein (immune saturation, local depletion, the ``all run out at
once'' intuition) are exactly the effects the branching-structure
comparison of \S\ref{sec:flooding} deliberately excludes, and should be
named wherever the flooding advantage is invoked.

\subsection{Open problems}
\label{sec:open-problems}

\begin{enumerate}
\item \textbf{A conceptual proof that burst size equals QSD.} The algebra
of \S\ref{sec:burst=qsd} is complete; a structural argument --- presumably
size-biasing composed with the burst intensity, or a Yaglom-type limit
\cite{yaglom1947certain} --- would turn a coincidence into a theorem one
can quote.
\item \textbf{Two-type killed second moments and the two-type $D$.} The
release kernel of a two-type renewal model needs
$\mean{X_a^2\mathbf1\{\tau_c>a\}}$, $\mean{X_aY_a\mathbf1\{\tau_c>a\}}$,
$\mean{Y_a^2\mathbf1\{\tau_c>a\}}$, and the analogue of $D(t)$
(extinction without catastrophe). The moment hierarchy does not close
upward, exactly as in one type; the one-type escape was to differentiate
the Riccati equation, and the two-type route should be to differentiate the
backward system of Chapter~5. This is the main technical obstacle between
this chapter and an eclipse-aware renewal BMVR, and it should not be
assumed easy.
\item \textbf{The flooding boundary against real parameters.} Map
$\delta(\beta+\mu)>\beta\mu$ for \yp{} in the macrophage, and compute the
\emph{size} of the advantage from \eqref{eq:flood}: does the pathogen sit
on the advantageous side, and by how much?
\item \textbf{Population-level variance.} Propagate $\mathrm{Var}(W)$ of
\eqref{eq:EW2} through the branching structure to a whole-host prediction:
for single-cell assays the first moments suffice, but for whole-host data
the variance is the main stochastic signal the BMVR cannot reproduce.
\item \textbf{Partial release and the flooding boundary in $q$.} Where on
the boolean continuum of \S\ref{sec:boolean} does the flooding advantage
survive? Does the closed-form structure persist at intermediate $q$?
\item \textbf{Logistic intracellular growth.} A referee will object that
linear birth lets the conditional load grow without bound. The defence is
that the burst hazard is proportional to load and so regulates it --- the
QSD result is precisely the statement that the conditional load stabilises
at $a/(a-1)$ --- but the sensitivity of $g(a)$ and $\Ihat(a)$ to a logistic
birth rate remains to be checked numerically.
\item \textbf{Literature positioning.} Verify whether Carruthers et al.\
\cite{carruthers2020stochastic} already couple an intracellular BDC to a
between-cell model --- the novelty claim of \S\ref{sec:renewal} depends on
the answer --- and position the renewal embedding against the
Pearson--Krapivsky--Perelson \cite{pearson2011stochastic},
Nelson--Gilchrist--Perelson, Gilchrist--Coombs
\cite{gilchrist2006evolution}, and Heffernan--Wahl lines of work, with the
BMVR's attribution settled between McLean, Nowak--Bangham, and Perelson et
al.
\end{enumerate}
